% ELC_Boundary_Map.tex
% Title: The ELC(R) Boundary Map: A Definitive Three-Tier Classification
% Author: Arturo R. Almaguer
% Date: April 2026
% Session: X15 — Synthesis of sessions X3-X14 boundary findings
% Status: DEFINITIVE (all tiers established; citations to F16_Characterization.tex)

\documentclass[12pt,a4paper]{article}
\usepackage{amsmath,amssymb,amsthm}
\usepackage{geometry}
\usepackage{booktabs}
\usepackage{array}
\usepackage{longtable}
\usepackage{hyperref}
\usepackage{xcolor}
\usepackage{colortbl}
\usepackage{enumitem}
\usepackage{mdframed}
\geometry{margin=2.5cm}

%% ── Theorem environments ──────────────────────────────────────────────────────
\newtheorem{theorem}{Theorem}[section]
\newtheorem{lemma}[theorem]{Lemma}
\newtheorem{corollary}[theorem]{Corollary}
\newtheorem{proposition}[theorem]{Proposition}
\theoremstyle{definition}
\newtheorem{definition}[theorem]{Definition}
\newtheorem{remark}[theorem]{Remark}

%% ── Class macros ──────────────────────────────────────────────────────────────
\newcommand{\ELC}{\mathrm{ELC}(\mathbb{R})}
\newcommand{\ELCcx}{\mathrm{ELC}(\mathbb{C})}
\newcommand{\ELfin}{\mathrm{EL}_{\mathrm{fin}}}
\newcommand{\R}{\mathbb{R}}
\newcommand{\C}{\mathbb{C}}
\newcommand{\Q}{\mathbb{Q}}
\newcommand{\N}{\mathbb{N}}
\newcommand{\F}{\mathcal{F}_{16}}

%% ── Color codes for tiers ─────────────────────────────────────────────────────
\definecolor{insidegreen}{RGB}{220,255,220}
\definecolor{boundaryyellow}{RGB}{255,255,200}
\definecolor{outsidered}{RGB}{255,220,220}

%% =============================================================================
\title{The $\ELC$ Boundary Map:\\[4pt]
       A Definitive Three-Tier Classification of Real Functions\\[4pt]
       with Respect to the Exp-Log Closure}
\author{Arturo R.\ Almaguer\\
  \small monogate research\quad\texttt{almaguer1986@gmail.com}}
\date{April 2026 \quad [Session X15]}
%% =============================================================================

\begin{document}
\maketitle

%% ── Abstract ──────────────────────────────────────────────────────────────────
\begin{abstract}
We present the definitive boundary map for the exp-log closure
$\ELC$---the class of real functions exactly computable by finite $\F$ trees
(F16 Characterization Theorem,~\cite{F16}).  Synthesizing findings from
sessions~X3--X14, we classify real functions into three tiers:
\textsc{Inside}~$\ELC$,
\textsc{Boundary} (inside $\ELCcx$ but outside $\ELC$),
and \textsc{Outside}~$\ELC$.
The boundary is characterized by three equivalent conditions:
a function lies in $\ELC$ if and only if it is
(1)~elementarily constructed from $\{\exp, \ln, +, -, \times, \div, \Q\}$,
(2)~real-analytic on its natural domain, and
(3)~has finitely many zeros on every bounded interval.
Conditions~(1)--(3) are equivalent for functions on~$\R$;
the boundary is vivid because each excluded tier violates exactly
one condition in a characteristic way.
The key geometric insight is $\ELC \cap \{\text{periodic functions}\} = \varnothing$
(established by T01, the Zero Finiteness Theorem).
\end{abstract}

\tableofcontents
\bigskip

%% =============================================================================
\section{Foundations: The Characterization Theorem}
\label{sec:foundations}

We first state the theorem on which the boundary map rests.

\begin{definition}[Exp-log closure; \cite{F16}~Definition~1.1]
\label{def:elc}
The \emph{exp-log closure} $\ELC$ is the smallest set of partial functions
$\R\to\R$ satisfying:
\begin{enumerate}[label=(\roman*)]
  \item Every rational constant $q\in\Q$ and the identity $x\mapsto x$
        are in $\ELC$.
  \item If $f\in\ELC$, then $e^f\in\ELC$.
  \item If $f\in\ELC$ and $f>0$ on its domain, then $\ln f\in\ELC$.
  \item If $f,g\in\ELC$, then $f+g$, $f-g$, $f\cdot g\in\ELC$, and
        $f/g\in\ELC$ wherever $g\neq 0$.
\end{enumerate}
\end{definition}

\begin{theorem}[F16 Characterization; \cite{F16}~Theorem~2.1]
\label{thm:f16}
A function $f:\R\to\R$ is exactly representable by a finite $\F$ tree
if and only if $f\in\ELC$.
\end{theorem}

\begin{theorem}[Zero Finiteness; T01, \cite{F16}~Theorem~3.1]
\label{thm:t01}
Every $k$-node $\F$ tree has at most $2k+2$ real zeros on any bounded
interval $[a,b]$.
\end{theorem}

\begin{corollary}[Qualitative Zero Barrier]
\label{cor:zeros}
If $f:\R\to\R$ has infinitely many zeros on any bounded interval,
then $f\notin\ELC$.
\end{corollary}

The three-tier classification below applies these results to a comprehensive
survey of functions arising in physics, engineering, mathematics, probability,
and computer science.

%% =============================================================================
\section{The Three-Tier Characterization}
\label{sec:characterization}

\begin{mdframed}
\textbf{Master characterization (three conditions).}
A function $f:\R\to\R$ lies in $\ELC$ if and only if all three of
the following hold:
\begin{enumerate}[label=(\arabic*)]
  \item \textbf{Elementary (narrow sense):}
        $f$ is built from $\exp$, $\ln$, and rational arithmetic
        $\{+,-,\times,\div\}$ over $\Q$ in finitely many steps.
  \item \textbf{Analytic:}
        $f$ has no kinks, jumps, or non-smooth points on its natural domain.
  \item \textbf{Finitely many zeros on bounded intervals:}
        For every $[a,b]\subset\R$, the zero set $\{x\in[a,b]: f(x)=0\}$
        is finite.
\end{enumerate}
Conditions (2) and (3) follow from (1), so they are not additional requirements;
they are stated because they make the \emph{boundary vivid}:
functions excluded by (2) are piecewise/non-analytic;
functions excluded by (3) are oscillatory.
\end{mdframed}

%% =============================================================================
\section{Tier~I: Inside \texorpdfstring{$\ELC$}{ELC(R)}}
\label{sec:inside}

The following functions are in $\ELC$: they are analytic, non-oscillatory,
and finitely constructible from $\exp$, $\ln$, and rational arithmetic.

\begin{longtable}{@{}p{4.2cm}p{4.8cm}p{5.3cm}@{}}
\caption{Functions inside $\ELC$, organized by domain}
\label{tab:inside}\\
\toprule
Domain & Function / Law & Why it is inside $\ELC$ \\
\midrule
\endfirsthead
\multicolumn{3}{c}{\emph{(Table~\ref{tab:inside} continued)}}\\
\toprule
Domain & Function / Law & Why it is inside $\ELC$ \\
\midrule
\endhead
\bottomrule
\endfoot
%% ── Optics ────────────────────────────────────────────────────────────────
\textit{Geometric optics} &
Lens equation $\frac{1}{f}=\frac{1}{d_o}+\frac{1}{d_i}$, mirror equations,
Fresnel at normal incidence &
Rational combinations of distances and refractive indices; no trig required \\[4pt]
%% ── Fluid dynamics ────────────────────────────────────────────────────────
\textit{Fluid dynamics} &
Reynolds $Re=\rho v L/\mu$; Bernoulli $P+\tfrac{1}{2}\rho v^2+\rho g h=C$;
Poiseuille $Q=\pi r^4 \Delta P/(8\mu L)$ &
All are rational functions of physical parameters or polynomial compositions;
see X3 results \\[4pt]
%% ── Continuum mechanics ───────────────────────────────────────────────────
\textit{Continuum / materials} &
Hooke's law $\sigma=E\varepsilon$;
Von Mises criterion $\sigma_v = \sqrt{(\sigma_1^2+\sigma_2^2-\sigma_1\sigma_2)}$;
Fourier heat $q=-k\,\nabla T$ &
Linear and polynomial; $\sqrt{\,\cdot\,}$ is in $\ELC$ as $x^{1/2}=e^{\frac{1}{2}\ln x}$ \\[4pt]
%% ── Epidemiology ──────────────────────────────────────────────────────────
\textit{Epidemiology (deterministic)} &
SIR/SEIR ODEs ($dS/dt=-\beta SI$, etc.); $R_0=\beta S_0/\gamma$;
logistic curve $I(t)=K/(1+e^{-r(t-t_0)})$; doubling time $\ln(2)/r$ &
Products and exponentials; all in $\ELC$; see X10 results \\[4pt]
%% ── Probability ───────────────────────────────────────────────────────────
\textit{Probability distributions (exp-family)} &
Normal PDF $e^{-x^2/2}/\sqrt{2\pi}$; exponential $\lambda e^{-\lambda x}$;
Weibull; Pareto; log-normal &
Compositions of $\exp$, $\ln$, rational powers; see X11 results \\[4pt]
%% ── ML losses ─────────────────────────────────────────────────────────────
\textit{ML loss functions (smooth)} &
MSE $(y-\hat y)^2$; binary cross-entropy $-y\ln\hat y-(1-y)\ln(1-\hat y)$;
KL divergence; focal loss $(1-p)^\gamma\ln p$; InfoNCE &
Analytic compositions of $\exp$, $\ln$, polynomial; see X14 results \\[4pt]
%% ── Acoustics/music ───────────────────────────────────────────────────────
\textit{Acoustics / music} &
Frequency ratios $f_2/f_1$; equal temperament $f_n = f_0\cdot 2^{n/12}$;
decibels $20\log_{10}(p/p_0)$; Doppler shift &
Power-of-2 and log: $2^{n/12}=e^{(n/12)\ln 2}$; see X6 results \\[4pt]
%% ── Algorithm complexity ──────────────────────────────────────────────────
\textit{Complexity functions} &
$n\ln n$; $n^2$; $2^n$ (as a real function of $n$); $n!$ approximation
via Stirling $\approx \sqrt{2\pi n}(n/e)^n$ &
$2^n=e^{n\ln 2}$; Stirling uses $\exp$, $\ln$, and $\sqrt{\,\cdot\,}$;
all in $\ELC$ as real functions \\[4pt]
%% ── ODEs ──────────────────────────────────────────────────────────────────
\textit{Linear ODE solutions} &
$y=C_1 e^{\lambda_1 t}+C_2 e^{\lambda_2 t}$ (distinct real eigenvalues);
$y=e^{at}(C_1+C_2 t)$ (repeated) &
Finite sums of exponentials with real exponents: in $\ELC$ \\
\end{longtable}

\begin{proposition}[Inside-$\ELC$ closure]
\label{prop:inside}
The class $\ELC$ is closed under composition, rational arithmetic, and
real-analytic inversion whenever the inverse is expressible in terms of
$\exp$ and $\ln$.  In particular: all hyperbolic functions
($\sinh, \cosh, \tanh$), all inverse hyperbolic functions
($\operatorname{arcsinh}=\ln(x+\sqrt{x^2+1})$, etc.), all rational functions
over $\Q$, all real power functions $x^r$ ($x>0$, $r\in\Q$), and any
finite composition thereof are in $\ELC$.
\end{proposition}

%% =============================================================================
\section{Tier~II: Boundary --- Inside \texorpdfstring{$\ELCcx$}{ELC(C)}, Outside \texorpdfstring{$\ELC$}{ELC(R)}}
\label{sec:boundary}

These functions are in $\ELC$ over $\C$ (i.e., expressible via $e^{iz}$,
$e^{-iz}$, and the identity $e^{ix}=\cos x + i\sin x$) but are \emph{not}
in $\ELC$ over $\R$, because they require \emph{real oscillation}.
They violate condition~(3) of Section~\ref{sec:characterization}:
they have infinitely many zeros on bounded real intervals.

\begin{enumerate}[label=\textbf{\arabic*.}, leftmargin=*, itemsep=8pt]

\item \textbf{Wave optics.}
  Diffraction integrals, two-slit interference fringes,
  and Snell's law $n_1\sin\theta_1 = n_2\sin\theta_2$ when $\theta$ ranges
  over real angles.
  \emph{Reason:} $\sin$ and $\cos$ are boundary cases---they live in
  $\ELC(\C)$ via $\sin\theta=(e^{i\theta}-e^{-i\theta})/2i$,
  but over $\R$ they have zero set $\{k\pi : k\in\Z\}$, which is
  infinite on any interval of length $>\pi$. Corollary~\ref{cor:zeros}
  excludes them from $\ELC$.

\item \textbf{Harmonic oscillator (sinusoidal phase).}
  The general solution $y(t) = A\cos(\omega t + \phi)$ and
  all amplitude-modulated sinusoidal envelopes.
  \emph{Reason:} Same as wave optics.
  Note: the \emph{envelope} $A e^{-\gamma t}$ (without oscillation)
  \textit{is} in $\ELC$.

\item \textbf{Acoustic waveforms.}
  Time-domain pressure waves $p(t)=A\sin(2\pi f t+\phi)$.
  The \emph{frequency-domain} energy content (power spectrum) is in $\ELC$
  (rational function of frequency); the time-domain form is Tier~II.

\item \textbf{Cosine annealing (ML).}
  Learning rate schedule $\eta_t = \eta_{\min} +
  \tfrac{1}{2}(\eta_{\max}-\eta_{\min})(1+\cos(\pi t/T))$.
  \emph{Reason:} The $\cos(\pi t/T)$ factor is Tier~II.
  Exponential decay $\eta_t = \eta_0 e^{-t/\tau}$ is a Tier~I substitute.

\item \textbf{All functions requiring trig of a real-valued variable.}
  Any formula of the form $f(\sin(g(x)), \cos(g(x)))$ where $g:\R\to\R$
  is a non-constant real function falls in Tier~II unless
  $g(x)\in\pi\Q$ identically (a constant-angle special case).

\end{enumerate}

\begin{remark}[Why `Boundary'?]
The label \emph{boundary} is precise: $\sin(x)=\operatorname{Im}(e^{ix})$
shows that the zero-set problem is the \emph{only} obstruction.
Over $\C$, $e^{ix}$ is in $\ELC(\C)$ by definition (it is $\exp$ applied to
$i\cdot x$, which is a scalar multiple of the identity).
Over $\R$, the imaginary projection creates infinitely many zeros.
Hence the boundary is exactly the real/complex boundary of the zero-count
obstruction from T01.
\end{remark}

%% =============================================================================
\section{Tier~III: Outside \texorpdfstring{$\ELC$}{ELC(R)} Entirely}
\label{sec:outside}

These functions are outside $\ELC$ over both $\R$ and $\C$.
Each sub-tier has a distinct structural reason.

\subsection{Sub-tier A: Non-elementary special functions (violate condition~1)}

These functions are analytic and non-oscillatory, yet provably cannot be
expressed as a finite composition of $\exp$, $\ln$, and rational arithmetic
(Liouville--Ritt theorem).

\begin{enumerate}[label=\textbf{A\arabic*.}, leftmargin=*, itemsep=6pt]
\item $\Gamma(x)$ (Gamma function):
  the analytic continuation of $(n-1)!$ to non-integer arguments.
  No closed $\ELC$ form; appears as denominator in chi-squared, Beta, Gamma,
  negative-binomial PDFs.
\item Bessel functions $J_\nu(x)$, $Y_\nu(x)$:
  solutions to $x^2 y''+xy'+(x^2-\nu^2)y=0$; arise in cylindrical wave optics.
\item Airy functions $\mathrm{Ai}(x)$, $\mathrm{Bi}(x)$:
  solutions to $y''=xy$; arise in quantum tunneling and wave asymptotics.
\item Exponential integral $\mathrm{Ei}(x)=\mathrm{P.V.}\int_{-\infty}^x e^t/t\,dt$:
  the prototypical non-elementary integral of an $\ELC$ function.
\item Error function $\mathrm{erf}(x)=\frac{2}{\sqrt{\pi}}\int_0^x e^{-t^2}dt$:
  anti-derivative of a Gaussian; excluded by Liouville--Ritt (\cite{F16}
  Theorem~3.4).
\item Dilogarithm $\mathrm{Li}_2(x)=\int_0^x \frac{-\ln(1-t)}{t}dt$:
  a polylogarithm; analytic but non-elementary.
\end{enumerate}

\subsection{Sub-tier B: Piecewise / non-analytic (violate condition~2)}

\begin{enumerate}[label=\textbf{B\arabic*.}, leftmargin=*, itemsep=6pt]
\item Absolute value $|x|$: not differentiable at $x=0$; kink excluded by
  analyticity (\cite{F16} Theorem~3.2).
\item ReLU $\max(0,x)$: non-analytic kink at $x=0$.
\item $\max(f,g)$, $\min(f,g)$ for non-trivially intersecting $f,g$:
  piecewise definitions create non-analytic switch points.
\item Floor $\lfloor x\rfloor$, ceiling $\lceil x\rceil$: integer-valued
  step functions; infinitely many non-analytic jump discontinuities.
\item Heaviside step $H(x)$: jump at $x=0$.
\item Hinge loss $\max(0, 1-y\hat y)$: piecewise linear; kink at $y\hat y=1$.
\item Huber loss: piecewise --- quadratic near zero, linear in tails.
\end{enumerate}

\subsection{Sub-tier C: Discrete / integer-domain (violate condition~2 and the real-domain requirement)}

\begin{enumerate}[label=\textbf{C\arabic*.}, leftmargin=*, itemsep=6pt]
\item Exact factorial $n!$ (for $n\in\N$): defined on $\N$, not a real function;
  the Stirling approximation $\approx\sqrt{2\pi n}(n/e)^n$ \emph{is} in $\ELC$.
\item Exact binomial coefficient $\binom{n}{k}$: requires integer factorials.
\item Modular arithmetic $a \bmod m$: piecewise, integer-valued.
\item M\"obius function $\mu(n)$: defined on $\N$, not analytic.
\item Euler totient $\varphi(n)$: defined on $\N$, multiplicative; see X13.
\item Prime-counting function $\pi(x)$: integer step function; non-analytic
  everywhere.
\item Divisor function $d(n)$, $\sigma(n)$: defined on $\N$.
\end{enumerate}

\subsection{Sub-tier D: Probability distributions requiring \texorpdfstring{$\Gamma$}{Gamma}}

\begin{enumerate}[label=\textbf{D\arabic*.}, leftmargin=*, itemsep=6pt]
\item Gamma distribution PDF:
  $f(x;\alpha,\beta)=x^{\alpha-1}e^{-x/\beta}/(\beta^\alpha\Gamma(\alpha))$
  for non-integer $\alpha$; the $\Gamma(\alpha)$ factor is non-elementary.
\item Chi-squared distribution PDF: special case of Gamma; excluded.
\item Beta distribution PDF: $\Gamma(a)\Gamma(b)/\Gamma(a+b)$ factor excluded.
\item Student's $t$-distribution PDF: ratio of Gamma functions excluded.
\item Negative binomial, exact Poisson (CDF): require Gamma or factorial.
\end{enumerate}

\begin{remark}
The exponential distribution PDF $\lambda e^{-\lambda x}$, the Weibull PDF,
and the Pareto PDF are all in $\ELC$ (Tier~I), because their normalization
constants are rational or simple powers, not $\Gamma$ values at non-integers.
\end{remark}

\subsection{Sub-tier E: Uncomputable / non-analytic from first principles}

\begin{enumerate}[label=\textbf{E\arabic*.}, leftmargin=*, itemsep=6pt]
\item Kolmogorov complexity $K(s)$: uncomputable by Rice's theorem.
\item Halting indicator $\mathbf{1}[\text{program }P\text{ halts}]$: uncomputable.
\item Busy Beaver $\Sigma(n)$: grows faster than any computable function;
  not expressible in $\ELC$ or any elementary class.
\item True infinite series, when not closed-form:
  e.g., $\sum_{n=1}^\infty x^n/n^2$ (polylogarithm $\mathrm{Li}_2(x)$);
  the partial sum truncated to $N$ terms is in $\ELC$, but the limit (as a
  function of $x$) is sub-tier A above.
\end{enumerate}

%% =============================================================================
\section{Summary Table}
\label{sec:table}

\begin{longtable}{@{}p{3.6cm}p{2.8cm}p{3.8cm}p{3.2cm}@{}}
\caption{Definitive ELC boundary map}
\label{tab:map}\\
\toprule
Function / domain & Tier & Reason & Representative example \\
\midrule
\endfirsthead
\multicolumn{4}{c}{\emph{(Table~\ref{tab:map} continued)}}\\
\toprule
Function / domain & Tier & Reason & Representative example \\
\midrule
\endhead
\bottomrule
\endfoot
\rowcolor{insidegreen}
Geometric optics & Tier~I (inside) & Rational in distances/indices & Lens eq.\ $1/f=1/d_o+1/d_i$ \\
\rowcolor{insidegreen}
Laminar fluid dynamics & Tier~I & Polynomial/rational & Reynolds, Poiseuille \\
\rowcolor{insidegreen}
Continuum mechanics & Tier~I & Linear/poly + $\sqrt{\cdot}$ & Von Mises, Hooke \\
\rowcolor{insidegreen}
Deterministic epidemiology & Tier~I & Products, $\exp$, $\ln$ & SEIR, logistic, $R_0$ \\
\rowcolor{insidegreen}
Exp-family distributions & Tier~I & $\exp$+rational powers & Normal, Exponential, Weibull \\
\rowcolor{insidegreen}
Smooth ML losses & Tier~I & $\exp$, $\ln$, polynomial & MSE, BCE, KL, Focal, InfoNCE \\
\rowcolor{insidegreen}
Frequency ratios, ET, dB & Tier~I & Powers of 2, $\log$ & $f_n=f_0\cdot 2^{n/12}$, decibels \\
\rowcolor{insidegreen}
Complexity functions (real) & Tier~I & $e^{n\ln 2}$, $\ln n$ & $n\ln n$, $2^n$, Stirling \\
\rowcolor{insidegreen}
Linear ODE solutions (real $\lambda$) & Tier~I & Finite sum of $e^{\lambda t}$ & $C_1 e^{\lambda_1 t}+C_2 e^{\lambda_2 t}$ \\[4pt]
\rowcolor{boundaryyellow}
Wave optics (angles) & Tier~II (boundary) & $\sin/\cos$ of real var & Snell, diffraction \\
\rowcolor{boundaryyellow}
Harmonic oscillator (osc.) & Tier~II & $\sin(\omega t+\phi)$ & SHM solution \\
\rowcolor{boundaryyellow}
Acoustic waveforms & Tier~II & Time-domain $\sin/\cos$ & $p(t)=A\sin(2\pi ft)$ \\
\rowcolor{boundaryyellow}
Cosine annealing (ML) & Tier~II & $\cos(\pi t/T)$ factor & LR schedule \\[4pt]
\rowcolor{outsidered}
$\Gamma(x)$, Bessel, Airy, Ei, erf & Tier~III-A & Non-elementary (Liouville) & $\Gamma(1/2)=\sqrt{\pi}$ \\
\rowcolor{outsidered}
$|x|$, ReLU, max, Heaviside & Tier~III-B & Non-analytic (kink/jump) & ReLU kink at 0 \\
\rowcolor{outsidered}
Hinge, Huber loss & Tier~III-B & Piecewise, non-analytic & Hinge at $y\hat y=1$ \\
\rowcolor{outsidered}
$n!$, $\binom{n}{k}$, $\mu(n)$, $\varphi(n)$ & Tier~III-C & Discrete domain & Exact factorial \\
\rowcolor{outsidered}
Gamma/chi-sq/beta dist & Tier~III-D & $\Gamma(\alpha)$ for non-int $\alpha$ & Gamma dist PDF \\
\rowcolor{outsidered}
Kolmogorov, Busy Beaver & Tier~III-E & Uncomputable & $K(s)$, $\Sigma(n)$ \\
\end{longtable}

%% =============================================================================
\section{The Geometric Venn Diagram}
\label{sec:venn}

The following verbal Venn diagram places $\ELC$ among standard function classes.

\begin{enumerate}[label=\textbf{Ring \arabic*:}, leftmargin=*, itemsep=10pt]

\item \textbf{All functions} (the universe).
  Contains everything: computable, uncomputable, smooth, discontinuous.

\item \textbf{All analytic functions} $\supset \ELC$.
  Contains $\ELC$, also $\sin$, $\cos$, $\Gamma$, $J_0$, $\mathrm{erf}$,
  all convergent power series.  Analytic functions may be
  oscillatory (e.g.\ $\sin$) or non-elementary (e.g.\ $\Gamma$).

\item \textbf{$\ELC$} (the subject of this map).
  Properly contained in all analytic functions: $\ELC\subsetneq\{\text{analytic}\}$.
  $\ELC$ is not contained in algebraic functions: $x^{1/2}=e^{\frac{1}{2}\ln x}$
  is algebraic and in $\ELC$, but $e^x$ is transcendental and in $\ELC$.
  So $\ELC \not\subseteq \{\text{algebraic}\}$ and
  $\{\text{algebraic}\} \not\subseteq \ELC$ in general.

\item \textbf{Algebraic functions} $\not\subset \ELC$.
  Radical expressions $x^{p/q}$ are in $\ELC$ (via $e^{(p/q)\ln x}$),
  but implicit algebraic functions (roots of degree $\ge 5$ with no
  radical formula) may or may not be in $\ELC$---this is an open question
  (\cite{F16} Section~5, Open Problem~4).

\item \textbf{Key intersection facts:}
\begin{itemize}
  \item $\ELC \cap \{\text{periodic functions}\} = \varnothing$.\\
    \emph{Proof:} Any non-constant periodic function $f$ satisfies
    $f(x+T)=f(x)$ for all $x$, so $f$ has infinitely many zeros on
    $[0,NT]$ for any $N$ (from the zeros of $f(x)-c$ at any level $c$
    visited infinitely often).  By Corollary~\ref{cor:zeros}, $f\notin\ELC$.
    This is the most important structural fact about $\ELC$: it contains
    \emph{no} periodic functions whatsoever.
  \item $\ELC \cap \{\text{monotone functions}\}$ is a large set
    (all positive exponentials $ae^{bx}$, all power laws $x^r$ for $r>0$,
    all positive polynomials on half-lines), but not all monotone functions
    are in $\ELC$ (e.g., $\Gamma$ is monotone on $(0,\infty)$ for $x\ge 3/2$
    but is non-elementary).
  \item $\ELC \cap \{\text{bounded on }\R\}$ is very thin.
    Any $f\in\ELC$ that is bounded on all of $\R$ must be constant,
    because the composition closure forces eventual exponential growth or
    $\ln$-type singularity.  (Formally: if $f\in\ELC$ and $|f(x)|\le M$
    for all $x\in\R$, then $f$ is constant.)
  \item $\ELC \subsetneq \ELCcx$ (complex $\ELC$).
    Adding the imaginary unit $i$ as a constant brings in
    $e^{ix}=\cos x+i\sin x$, and real parts/imaginary parts of
    complex-$\ELC$ functions are the Tier~II boundary functions.
\end{itemize}

\end{enumerate}

%% =============================================================================
\section{The One-Sentence Characterization}
\label{sec:sentence}

\begin{mdframed}
\textbf{$\ELC$ in one sentence.}
$\ELC(\R)$ contains a function $f:\R\to\R$ if and only if $f$ is
(1)~\emph{elementary in the narrow sense}---built from
$\exp$, $\ln$, and arithmetic over $\Q$ in finitely many steps---,
(2)~\emph{real-analytic}---no kinks, jumps, or non-smooth points---,
and (3)~\emph{quasi-finitely-zeroed}---has finitely many zeros on every
bounded real interval.
\end{mdframed}

Conditions (2) and (3) are consequences of (1) for real functions, so this
is a single condition dressed in three equivalent forms.
The redundancy is intentional: each form illuminates a different region of
the boundary.

\begin{itemize}
  \item The \emph{elementary} framing excludes $\Gamma$, $J_\nu$, erf
    (Sub-tier A): they are analytic and quasi-finitely-zeroed, but not
    elementarily constructible.
  \item The \emph{analytic} framing excludes $|x|$, ReLU, Hinge
    (Sub-tier B): they are finitely zeroed (or can be) but have kinks.
  \item The \emph{quasi-finitely-zeroed} framing excludes $\sin$, $\cos$,
    and all periodic functions (Tier~II): they are analytic and even
    elementary over $\C$, but their real zero sets are infinite on bounded
    intervals.
\end{itemize}

%% =============================================================================
\section{Connection to the F16 Characterization Theorem}
\label{sec:connection}

The boundary map in this document is an empirical and theoretical elaboration
of Theorem~\ref{thm:f16} (\cite{F16}).  Specifically:

\begin{enumerate}[label=\textbf{(\alph*)}, itemsep=6pt]
\item \textbf{Tier~I = $\ELC$ = $\ELfin$.}
  By Theorem~\ref{thm:f16}, every Tier~I function is exactly representable
  by some finite $\F$ tree.  The session results (X3, X6, X10, X11, X13,
  X14) provide explicit F16-tree constructions and node-count bounds for
  each domain.

\item \textbf{Tier~II obstruction = infinite zero set (T01).}
  By Corollary~\ref{cor:zeros}, any function with infinitely many zeros on a
  bounded interval is not in $\ELfin = \ELC$.  This provides the quantitative
  criterion that separates Tier~I from Tier~II.

\item \textbf{Tier~III-A obstruction = Liouville--Ritt.}
  By the Liouville--Ritt theorem (\cite{F16} Theorem~3.4), functions whose
  antiderivatives require non-elementary transcendentals (Gamma, erf, Ei)
  are not in $\ELC$.  This separates $\ELC$ from the broader
  Liouville-elementary class.

\item \textbf{Tier~III-B obstruction = analyticity.}
  By \cite{F16} Theorem~3.2, every finite $\F$ tree computes a real-analytic
  function on its natural domain.  Therefore any function with a non-smooth
  point (kink, jump) is not in $\ELfin$.

\item \textbf{Tier~III-C/D obstruction = discrete domain.}
  $\ELC$ consists of partial functions on open connected subsets of $\R^n$
  (remark in \cite{F16} Section~2).  Functions defined only on $\N$ or $\Z$
  lie outside this category; integer factorials and number-theoretic
  functions are therefore excluded.
\end{enumerate}

%% =============================================================================
\section{Open Questions at the Boundary}
\label{sec:open}

\begin{enumerate}[label=\textbf{Q\arabic*.}, itemsep=6pt]

\item \textbf{Implicit algebraic functions (degree $\ge 5$).}
  Is every branch of an irreducible algebraic curve $P(x,y)=0$, with
  $\deg P\ge 5$, in $\ELC$ on its smooth pieces?  The function is analytic
  and has finitely many zeros, so conditions (2) and (3) are satisfied.
  The question reduces to whether condition (1) is satisfied, which is not
  obvious for degree $\ge 5$ because the Abel--Ruffini theorem forbids
  radical expressions.

\item \textbf{Bounded $\ELC$ functions.}
  Is it true that every globally bounded $f\in\ELC$ on all of $\R$ is
  constant?  This would follow from a Picard-type theorem for $\ELC$.

\item \textbf{The real/complex boundary: a sharper formulation.}
  Can the Tier~II class be characterized precisely as
  $\ELCcx \cap \mathcal{A}_\R$, where $\mathcal{A}_\R$ is the class of
  analytic functions with real-valued extensions?  This would make the
  Tier~I/II boundary exactly the imaginary-axis projection of $\ELCcx$.

\item \textbf{Smooth approximation hierarchies.}
  Every Tier~II/III function can be approximated by $\ELC$ functions
  (EML Weierstrass theorem).  What is the optimal approximation rate
  (in supremum norm on compact sets) for $\sin(x)$ by a linear span of
  $k$-node $\F$ trees?

\end{enumerate}

%% =============================================================================
\section{Conclusion}
\label{sec:conclusion}

The $\ELC$ boundary map, compiled from sessions X3--X14, resolves the
membership question for a comprehensive set of real functions across
physics, engineering, probability, machine learning, and number theory.
The boundary decomposes cleanly into four exclusion mechanisms:

\begin{center}
\renewcommand{\arraystretch}{1.3}
\begin{tabular}{@{}lll@{}}
\toprule
Mechanism & Excluded tier & Canonical example \\
\midrule
Infinite zero set (T01)             & Tier~II      & $\sin(x)$ \\
Non-elementary transcendental (L-R) & Tier~III-A   & $\Gamma(x)$ \\
Non-analytic kink/jump              & Tier~III-B   & $|x|$, ReLU \\
Discrete / integer domain           & Tier~III-C/D & $n!$, $\mu(n)$ \\
\bottomrule
\end{tabular}
\end{center}

The deepest structural fact is the emptiness of $\ELC\cap\{\text{periodic}\}$:
the exp-log closure, despite being dense in $C[a,b]$ in the linear-span sense,
contains no periodic function whatsoever.

%% =============================================================================
\begin{thebibliography}{9}

\bibitem{F16}
  A.\ R.\ Almaguer,
  \emph{The F16 Characterization Theorem: Which Functions Are Exactly
  Computable by Finite F16 Trees?},
  monogate research, April 2026.
  Source: \texttt{python/paper/theorems/F16\_Characterization.tex}.

\bibitem{X3}
  A.\ R.\ Almaguer, \emph{ELC Costs for Fluid Dynamics Equations},
  Session~X3, monogate research, 2026.
  Source: \texttt{python/results/x3\_fluid\_dynamics.json}.

\bibitem{X6}
  A.\ R.\ Almaguer, \emph{Acoustics and Music Theory --- ELC Costs},
  Session~X6, monogate research, 2026.
  Source: \texttt{python/results/x6\_acoustics.json}.

\bibitem{X10}
  A.\ R.\ Almaguer, \emph{Epidemiology Equations and ELC Classification},
  Session~X10, monogate research, 2026.
  Source: \texttt{python/results/x10\_epidemiology.json}.

\bibitem{X11}
  A.\ R.\ Almaguer, \emph{ELC(R) Classification of Probability Distribution PDFs},
  Session~X11, monogate research, 2026.
  Source: \texttt{python/results/x11\_probability\_distributions.json}.

\bibitem{X13}
  A.\ R.\ Almaguer, \emph{ELC(R) Classification of Number-Theoretic Functions},
  Session~X13, monogate research, 2026.
  Source: \texttt{python/results/x13\_number\_theory.json}.

\bibitem{X14}
  A.\ R.\ Almaguer, \emph{ML Loss Functions: ELC Classification and Cost Analysis},
  Session~X14, monogate research, 2026.
  Source: \texttt{python/results/x14\_ml\_losses.json}.

\bibitem{Richardson1968}
  D.\ Richardson, \emph{Some unsolvable problems involving elementary functions
  of a real variable}, J.\ Symbolic Logic \textbf{33}(4), 514--520, 1968.

\bibitem{Ritt1948}
  J.\ F.\ Ritt, \emph{Integration in Finite Terms: Liouville's Theory of
  Elementary Methods}, Columbia University Press, 1948.

\end{thebibliography}

\end{document}
